\documentclass[11pt,a4paper]{article}
% ============================================================
%  Basado en la plantilla de las Prácticas 8-11 (mismos colores/estilos)
% ============================================================

% ---------- Idioma y codificación ----------
\usepackage[utf8]{inputenc}
\usepackage[T1]{fontenc}
\usepackage[spanish,es-noshorthands,es-tabla]{babel}

% ---------- Tipografía y layout ----------
\usepackage{lmodern}
\usepackage[sfdefault,light]{roboto}
\usepackage{geometry}
\geometry{a4paper, top=3.1cm, bottom=2.8cm, left=2.4cm, right=2.4cm, headheight=1.6cm}
\usepackage{ragged2e}
\usepackage{enumitem}
\usepackage{graphicx}
\usepackage{float}
\usepackage{booktabs}
\usepackage{array}
\usepackage[most]{tcolorbox}
\usepackage{tikz}
\usetikzlibrary{shadows,calc}
\usepackage{titlesec}
\usepackage{fancyhdr}
\usepackage{fontawesome5}
\usepackage[colorlinks=true,breaklinks=true]{hyperref}
\usepackage{caption}
\usepackage{microtype}
\usepackage{xurl}
\usepackage{listings}
\usepackage{amsmath,amssymb}
\usepackage{needspace}
\usepackage{icomma}

% ---------- Paleta de colores (tema "coolnight" de tu terminal Kitty) ----------
\definecolor{cNavy}{HTML}{011423}
\definecolor{cTealDeep}{HTML}{064B5A}
\definecolor{cTeal}{HTML}{097289}
\definecolor{cBlue}{HTML}{0FC5ED}
\definecolor{cCyan}{HTML}{0FC5ED}
\definecolor{cMint}{HTML}{44FFB1}
\definecolor{cPurple}{HTML}{A277FF}
\definecolor{cRed}{HTML}{E52E2E}
\definecolor{cBg}{HTML}{EAF6F8}
\definecolor{cText}{HTML}{01111E}
\definecolor{cOutBg}{HTML}{F4F4F4}

\hypersetup{
  linkcolor=cTealDeep,
  citecolor=cTealDeep,
  urlcolor=cBlue,
  pdftitle={Trabajo final: regresión lineal y logística, componentes principales y conglomerados},
  pdfauthor={Miguel Deza}
}

\color{cText}
\emergencystretch=2.5em

% ---------- Encabezado / pie de página ----------
\pagestyle{fancy}
\fancyhf{}
\renewcommand{\headrulewidth}{0pt}
\renewcommand{\footrulewidth}{0pt}
\fancyhead[L]{\small\color{cTealDeep}\textbf{\faGraduationCap\ Maestría en Ciencia de Datos}}
\fancyhead[R]{\small\color{cTeal}Estadística en Ciencia de Datos -- Trabajo final}
\fancyfoot[C]{%
  \begin{tikzpicture}[remember picture]
    \node[circle, fill=cTealDeep, minimum size=1.7em, inner sep=0pt] (pg) {\color{white}\small\thepage};
  \end{tikzpicture}%
}
\renewcommand{\headrule}{{\color{cCyan}\hrule height 1.4pt}\vspace{1pt}{\color{cBlue}\hrule height 0.5pt}}

% ---------- Estilo de secciones ----------
\titleformat{\section}
  {\normalfont\Large\bfseries\color{cTealDeep}}
  {\tikz[baseline=-0.6ex]{\node[fill=cCyan, text=white, rounded corners=1.5pt, inner sep=3pt, font=\bfseries]{\thesection};}}
  {0.6em}{}
  [{\vspace{-0.3em}\color{cCyan}\titlerule[1.2pt]\vspace{0.3em}}]

\titleformat{\subsection}
  {\normalfont\large\bfseries\color{cTeal}}
  {\color{cTeal}\thesubsection}{0.6em}{}

\titlespacing*{\section}{0pt}{1.4em}{0.8em}
\titlespacing*{\subsection}{0pt}{1.1em}{0.5em}

\captionsetup{labelfont={bf,color=cTealDeep}, font=small, textfont=it}

% ---------- Bloques de código R y salida de consola ----------
\lstdefinestyle{styleR}{
  language=R,
  backgroundcolor=\color{cBg},
  basicstyle=\ttfamily\footnotesize\color{cText},
  keywordstyle=\color{cTealDeep}\bfseries,
  commentstyle=\color{gray}\itshape,
  stringstyle=\color{cTeal},
  numbers=none,
  breaklines=true,
  frame=l,
  framerule=2.2pt,
  rulecolor=\color{cCyan},
  xleftmargin=8pt,
  framexleftmargin=6pt,
  aboveskip=6pt, belowskip=6pt,
  showstringspaces=false,
  columns=fullflexible, keepspaces=true
}

\lstdefinestyle{styleOut}{
  backgroundcolor=\color{cOutBg},
  basicstyle=\ttfamily\scriptsize\color{cText},
  breaklines=true,
  frame=l,
  framerule=2.2pt,
  rulecolor=\color{cTealDeep},
  xleftmargin=8pt,
  framexleftmargin=6pt,
  aboveskip=2pt, belowskip=6pt,
  showstringspaces=false,
  columns=fullflexible, keepspaces=true
}

% ---------- Preguntas ----------
\setlist{itemsep=1pt, topsep=3pt}
\graphicspath{{img/}}
\setlength{\parskip}{3pt}

\newtcolorbox{enunciado}[1][]{enhanced, breakable, colback=cBg, colframe=cTealDeep,
  boxrule=0pt, leftrule=3pt, arc=2pt, left=8pt, right=8pt, top=5pt, bottom=5pt,
  fonttitle=\bfseries, coltitle=white, colbacktitle=cTealDeep,
  title={Pregunta #1}}
\newtcolorbox{resultado}{enhanced, breakable, colback=cMint!12, colframe=cTealDeep,
  boxrule=0.5pt, arc=2pt, left=8pt, right=8pt, top=4pt, bottom=4pt}

% \pregunta{etiqueta}{texto del índice}{enunciado}
\newcommand{\pregunta}[3]{%
  \Needspace{12\baselineskip}\medskip
  \phantomsection
  \addcontentsline{toc}{subsection}{#1 \ #2}%
  \begin{enunciado}[#1]#3\end{enunciado}\noindent}

% portada de sección: página completa con fondo oscuro
\newcommand{\seccion}[4]{%
  \clearpage\newgeometry{margin=0cm}
  \refstepcounter{section}\phantomsection
  \addcontentsline{toc}{section}{\protect\numberline{\thesection}\protect\faIcon{#1}\ #2}%
  \thispagestyle{empty}\noindent
  \begin{tikzpicture}[remember picture, overlay]
    \fill[cNavy] (current page.south west) rectangle (current page.north east);
    \fill[cTealDeep, opacity=0.55] (current page.north west) ++(1.5,-1) circle (4.2);
    \fill[cCyan, opacity=0.14] (current page.north west) ++(4.2,-0.4) circle (2.2);
    \fill[cPurple, opacity=0.26] (current page.south east) ++(-1.2,1.4) circle (3.6);
    \fill[cMint, opacity=0.13] (current page.south east) ++(-3.6,0.3) circle (2.0);
    \draw[cCyan, opacity=0.25, line width=0.8pt] (current page.center) ++(5.2,4.6) circle (1.1);
    \node[anchor=center] at ($(current page.center)+(0,5.2)$)
      {\color{cTealDeep!80!cCyan}\bfseries\fontsize{150}{150}\selectfont \thesection};
    \node[circle, fill=cTealDeep, draw=cCyan, line width=1.3pt, minimum size=2.4cm] (ic)
      at ($(current page.center)+(0,1.6)$) {};
    \node at (ic) {\color{cMint}\fontsize{30}{30}\selectfont\faIcon{#1}};
    \node[anchor=center] at ($(current page.center)+(0,-0.4)$)
      {\color{cCyan}\bfseries\fontsize{11}{13}\selectfont S E C C I Ó N};
    \node[anchor=north, text width=14cm, align=center] (tit) at ($(current page.center)+(0,-0.9)$)
      {\hyphenpenalty=10000 {\color{white}\bfseries\fontsize{30}{36}\selectfont #3}};
    \draw[cCyan, line width=1.4pt] ($(tit.south)+(-3.2,-0.3)$) -- ++(6.4,0);
    \node[anchor=north, text width=14cm, align=center] at ($(tit.south)+(0,-0.6)$)
      {\color{cMint}\fontsize{13}{16}\selectfont #4};
  \end{tikzpicture}%
  \clearpage\restoregeometry}

% código y salidas: las líneas de cada bloque se leen de rangos.tex,
% generado por generar-salidas.R desde RLM.R, Logistica.R, CP.R y Conglomerados.R
\expandafter\def\csname r@rlm-datos\endcsname{28-34}
\expandafter\def\csname s@rlm-datos\endcsname{RLM.R}
\expandafter\def\csname r@rlm-ajuste\endcsname{38-41}
\expandafter\def\csname s@rlm-ajuste\endcsname{RLM.R}
\expandafter\def\csname r@rlm-r2\endcsname{45-47}
\expandafter\def\csname s@rlm-r2\endcsname{RLM.R}
\expandafter\def\csname r@rlm-supuestos\endcsname{51-53}
\expandafter\def\csname s@rlm-supuestos\endcsname{RLM.R}
\expandafter\def\csname r@rlm-diag1\endcsname{57-65}
\expandafter\def\csname s@rlm-diag1\endcsname{RLM.R}
\expandafter\def\csname r@rlm-diag2\endcsname{69-77}
\expandafter\def\csname s@rlm-diag2\endcsname{RLM.R}
\expandafter\def\csname r@rlm-vif\endcsname{81-81}
\expandafter\def\csname s@rlm-vif\endcsname{RLM.R}
\expandafter\def\csname r@rlm-backward\endcsname{85-85}
\expandafter\def\csname s@rlm-backward\endcsname{RLM.R}
\expandafter\def\csname r@rlm-final\endcsname{89-93}
\expandafter\def\csname s@rlm-final\endcsname{RLM.R}
\expandafter\def\csname r@rlm-final-sup\endcsname{97-100}
\expandafter\def\csname s@rlm-final-sup\endcsname{RLM.R}
\expandafter\def\csname r@rlm-diag-final\endcsname{104-112}
\expandafter\def\csname s@rlm-diag-final\endcsname{RLM.R}
\expandafter\def\csname r@rlm-sensibilidad\endcsname{116-121}
\expandafter\def\csname s@rlm-sensibilidad\endcsname{RLM.R}
\expandafter\def\csname r@log-datos\endcsname{28-33}
\expandafter\def\csname s@log-datos\endcsname{Logistica.R}
\expandafter\def\csname r@log-justifica\endcsname{37-46}
\expandafter\def\csname s@log-justifica\endcsname{Logistica.R}
\expandafter\def\csname r@log-ajuste\endcsname{50-53}
\expandafter\def\csname s@log-ajuste\endcsname{Logistica.R}
\expandafter\def\csname r@log-or\endcsname{57-58}
\expandafter\def\csname s@log-or\endcsname{Logistica.R}
\expandafter\def\csname r@log-lr\endcsname{62-64}
\expandafter\def\csname s@log-lr\endcsname{Logistica.R}
\expandafter\def\csname r@log-confusion\endcsname{68-74}
\expandafter\def\csname s@log-confusion\endcsname{Logistica.R}
\expandafter\def\csname r@log-roc\endcsname{78-82}
\expandafter\def\csname s@log-roc\endcsname{Logistica.R}
\expandafter\def\csname r@log-roc-fig\endcsname{86-98}
\expandafter\def\csname s@log-roc-fig\endcsname{Logistica.R}
\expandafter\def\csname r@log-cv\endcsname{102-111}
\expandafter\def\csname s@log-cv\endcsname{Logistica.R}
\expandafter\def\csname r@acp-datos\endcsname{28-34}
\expandafter\def\csname s@acp-datos\endcsname{CP.R}
\expandafter\def\csname r@acp-cor\endcsname{38-43}
\expandafter\def\csname s@acp-cor\endcsname{CP.R}
\expandafter\def\csname r@acp-corplot\endcsname{47-51}
\expandafter\def\csname s@acp-corplot\endcsname{CP.R}
\expandafter\def\csname r@acp-pca\endcsname{55-61}
\expandafter\def\csname s@acp-pca\endcsname{CP.R}
\expandafter\def\csname r@acp-scree\endcsname{65-76}
\expandafter\def\csname s@acp-scree\endcsname{CP.R}
\expandafter\def\csname r@acp-cargas\endcsname{80-81}
\expandafter\def\csname s@acp-cargas\endcsname{CP.R}
\expandafter\def\csname r@acp-cargas-fig\endcsname{85-94}
\expandafter\def\csname s@acp-cargas-fig\endcsname{CP.R}
\expandafter\def\csname r@acp-individuos\endcsname{98-114}
\expandafter\def\csname s@acp-individuos\endcsname{CP.R}
\expandafter\def\csname r@acp-biplot\endcsname{118-121}
\expandafter\def\csname s@acp-biplot\endcsname{CP.R}
\expandafter\def\csname r@acp-regresion\endcsname{125-131}
\expandafter\def\csname s@acp-regresion\endcsname{CP.R}
\expandafter\def\csname r@clu-datos\endcsname{28-37}
\expandafter\def\csname s@clu-datos\endcsname{Conglomerados.R}
\expandafter\def\csname r@clu-dist\endcsname{41-44}
\expandafter\def\csname s@clu-dist\endcsname{Conglomerados.R}
\expandafter\def\csname r@clu-dist-fig\endcsname{48-59}
\expandafter\def\csname s@clu-dist-fig\endcsname{Conglomerados.R}
\expandafter\def\csname r@clu-ward\endcsname{63-68}
\expandafter\def\csname s@clu-ward\endcsname{Conglomerados.R}
\expandafter\def\csname r@clu-dendro\endcsname{72-81}
\expandafter\def\csname s@clu-dendro\endcsname{Conglomerados.R}
\expandafter\def\csname r@clu-perfil-hc\endcsname{85-94}
\expandafter\def\csname s@clu-perfil-hc\endcsname{Conglomerados.R}
\expandafter\def\csname r@clu-kmeans\endcsname{98-109}
\expandafter\def\csname s@clu-kmeans\endcsname{Conglomerados.R}
\expandafter\def\csname r@clu-centroides\endcsname{113-121}
\expandafter\def\csname s@clu-centroides\endcsname{Conglomerados.R}
\expandafter\def\csname r@clu-compara\endcsname{125-136}
\expandafter\def\csname s@clu-compara\endcsname{Conglomerados.R}
\expandafter\def\csname r@clu-plano\endcsname{140-151}
\expandafter\def\csname s@clu-plano\endcsname{Conglomerados.R}
\expandafter\def\csname r@clu-enfermedad\endcsname{155-161}
\expandafter\def\csname s@clu-enfermedad\endcsname{Conglomerados.R}
\expandafter\def\csname r@clu-enf-fig\endcsname{165-173}
\expandafter\def\csname s@clu-enf-fig\endcsname{Conglomerados.R}

\newcommand{\codigo}[1]{%
  \edef\figopts{linerange={\csname r@#1\endcsname}}\edef\figfile{\csname s@#1\endcsname}%
  \expandafter\lstinputlisting\expandafter[\figopts, style=styleR]{\figfile}}
\newcommand{\salida}[1]{\lstinputlisting[style=styleOut]{salidas/#1.txt}}
% figura a la izquierda y su código R a la derecha
\newcommand{\figcod}[2]{%
  \par\medskip\noindent
  \begin{minipage}[c]{0.47\linewidth}\centering
    \includegraphics[width=\linewidth]{#1}\captionof{figure}{#2}
  \end{minipage}\hfill
  \begin{minipage}[c]{0.5\linewidth}
    \edef\figopts{linerange={\csname r@#1\endcsname}}\edef\figfile{\csname s@#1\endcsname}%
    \expandafter\lstinputlisting\expandafter[\figopts, style=styleR, basicstyle=\ttfamily\scriptsize\color{cText}]{\figfile}
  \end{minipage}\par\medskip}
\newcommand{\res}[1]{\begin{resultado}\textbf{\color{cTealDeep}Resultado.} #1\end{resultado}}
\newcommand{\var}[1]{\texttt{\detokenize{#1}}}

% ============================================================
% portada: \portada{icono}{título}{subtítulo}{descripción}
\newcommand{\portada}[4]{%
\begin{titlepage}
\newgeometry{margin=0cm}
\begin{tikzpicture}[remember picture, overlay]
  \fill[cNavy] (current page.south west) rectangle (current page.north east);
  \fill[cPurple, opacity=0.28] (current page.north east) ++(-1,1) circle (3.2);
  \fill[cCyan, opacity=0.22] (current page.north east) ++(-2.6,0.4) circle (2.1);
  \fill[cTealDeep, opacity=0.55] (current page.south west) ++(1,-1) circle (4.0);
  \fill[cMint, opacity=0.16] (current page.south west) ++(2.4,-0.2) circle (2.3);

  \node[circle, fill=cTealDeep, draw=cCyan, line width=1.3pt, minimum size=2.5cm] (logo) at ($(current page.north)+(0,-2.9)$) {};
  \node at (logo) {\color{cMint}\fontsize{28}{28}\selectfont\faIcon{#1}};

  \node[anchor=north, text width=14cm, align=center] at ($(current page.north)+(0,-4.7)$) {
    {\color{cCyan}\bfseries\fontsize{12}{14}\selectfont \faDatabase\ \ MAESTRÍA EN CIENCIA DE DATOS}
  };

  \node[anchor=center, text width=15cm, align=center] at ($(current page.center)+(0,1.2)$) {
    {\color{white}\bfseries\fontsize{30}{38}\selectfont #2}\\[0.35em]
    {\color{cMint}\bfseries\fontsize{17}{21}\selectfont \faIcon{#1}\ #3}
  };

  \draw[cCyan, line width=1.4pt] ($(current page.center)+(-3.6,-0.9)$) -- ++(7.2,0);

  \node[anchor=center, text width=13cm, align=center] at ($(current page.center)+(0,-1.9)$) {
    {\color{white!85!black}\fontsize{11}{14}\selectfont
     #4}
  };

  \node[anchor=south, text width=14cm, align=center] at ($(current page.south)+(0,2.6)$) {
    \begin{tabular}{r l}
      \color{cCyan}\bfseries \faUser\ Elaborado por: & \color{white} Miguel Deza \\[3pt]
      \color{cCyan}\bfseries \faGraduationCap\ Programa: & \color{white} Maestría en Ciencia de Datos -- UNSA \\[3pt]
      \color{cCyan}\bfseries \faBook\ Curso: & \color{white} Estadística en Ciencia de Datos \\[3pt]
      \color{cCyan}\bfseries \faChalkboardTeacher\ Docente: & \color{white} Carlos Antonio Abanto Valle \\[3pt]
      \color{cCyan}\bfseries \faIcon{calendar-alt}\ Fecha: & \color{white} 7 de noviembre de 2026 \\
    \end{tabular}
  };
\end{tikzpicture}
\end{titlepage}
\restoregeometry
}

\hypersetup{pdftitle={Regresión logística binaria}}
\begin{document}
\portada{heartbeat}{Regresión Logística\\ Binaria}{Odds ratios, matriz de confusión y curva ROC}{Base de datos: Heart Disease (Cleveland, UCI)\\ Máxima verosimilitud, razón de verosimilitudes, ROC/AUC}

\setcounter{tocdepth}{2}
\tableofcontents

\vspace{1.2em}
\noindent Todo el análisis está en \var{2_Regresion_Logistica_Binaria.R} (R 4.6.1; paquetes \var{lmtest} y \var{pROC}), que lee \var{heart.csv} (297 pacientes, sin valores faltantes). En cada pregunta muestro el bloque de código, su salida de consola y la interpretación.

\clearpage
\section*{Datos y modelo}
\[
\log\frac{P(Y_i=1\mid x_i)}{1-P(Y_i=1\mid x_i)}=\beta_0+\beta_1\,\text{edad}_i+\beta_2\,\text{sexo}_i
+\beta_3\,\text{colesterol}_i+\beta_4\,\text{presion\_reposo}_i
\]
\[
+\,\beta_5\,\text{frecuencia\_cardiaca\_max}_i+\beta_6\,\text{angina\_ejercicio}_i+\beta_7\,\text{depresion\_st}_i
\]
con $Y_i=1$ si el paciente tiene enfermedad cardíaca. \var{sexo} (referencia: mujer) y
\var{angina_ejercicio} (referencia: no) entran como factores.
\codigo{log-datos}
\salida{log-datos}
Las clases están casi balanceadas: 137 enfermos (46,1\,\%) y 160 sanos.

\paragraph{Justificación de las variables.} Antes de ajustar comparé cada covariable entre enfermos y sanos:
\codigo{log-justifica}
\salida{log-justifica}
Seis de las siete variables difieren claramente entre grupos: los enfermos son mayores, con más presión,
menor frecuencia cardíaca máxima, mayor depresión ST, más hombres (56\,\% vs.\ 26\,\% de enfermos) y más
angina de esfuerzo (76\,\% vs.\ 32\,\%). El colesterol no difiere de forma marginal ($p=0,17$), pero lo
mantengo porque es un factor de riesgo clásico y puede aportar ajustado por las demás. Todas son medidas
clínicas baratas, disponibles antes de un cateterismo.

\pregunta{a)}{Ajuste por máxima verosimilitud}{Ajuste el modelo de regresión logística binaria mediante el
método de máxima verosimilitud.}
\[
\ln L(\boldsymbol\beta)=\sum_{i=1}^{n}\Big[y_i\,x_i^\top\boldsymbol\beta-\ln\!\big(1+e^{x_i^\top\boldsymbol\beta}\big)\Big],
\qquad
\frac{\partial\ln L}{\partial\boldsymbol\beta}=\sum_{i=1}^{n}x_i\,(y_i-\pi_i)=\mathbf 0
\]
No tiene solución cerrada: \var{glm()} la resuelve por Fisher scoring (IRLS), que convergió en 5 iteraciones.
\codigo{log-ajuste}
\salida{log-ajuste}

\res{La devianza baja de 409,95 (modelo nulo) a 274,16 con 7 parámetros adicionales; AIC $=290,2$.}

\pregunta{b)}{Tabla de estimación y significancia}{Presente la tabla de estimación de los parámetros e
identifique las variables estadísticamente significativas al nivel $\alpha=0,05$.}
\begin{center}\small
\begin{tabular}{lrrrrc}
\toprule
Variable & $\hat\beta$ & EE & $z$ & $p$ & ¿Signif.? \\ \midrule
Intercepto & $-2,383$ & 2,065 & $-1,15$ & 0,249 & no \\
edad & 0,0199 & 0,0195 & 1,02 & 0,309 & no \\
sexo (hombre) & 1,652 & 0,363 & 4,55 & $<0,001$ & \textbf{sí} \\
colesterol & 0,0062 & 0,0030 & 2,04 & 0,041 & \textbf{sí} \\
presion\_reposo & 0,0142 & 0,0091 & 1,56 & 0,118 & no \\
frecuencia\_cardiaca\_max & $-0,0296$ & 0,0084 & $-3,54$ & $<0,001$ & \textbf{sí} \\
angina\_ejercicio (sí) & 1,325 & 0,337 & 3,93 & $<0,001$ & \textbf{sí} \\
depresion\_st & 0,583 & 0,148 & 3,94 & $<0,001$ & \textbf{sí} \\ \bottomrule
\end{tabular}
\end{center}
Edad y presión en reposo sí diferían entre grupos por separado, pero dejan de ser significativas al
ajustar: su información ya la cargan la frecuencia cardíaca máxima (que baja con la edad) y la depresión ST.

\res{Significativas al 5\,\%: \textbf{sexo, colesterol, frecuencia\_cardiaca\_max, angina\_ejercicio y
depresion\_st}.}

\pregunta{c)}{Interpretación de edad, angina\_ejercicio y depresion\_st}{Interprete los coeficientes asociados
a las variables edad, angina\_ejercicio y depresion\_st en términos de odds ratios.}
\[
OR_j=e^{\hat\beta_j}=\frac{\text{odds}(x_j+1)}{\text{odds}(x_j)}
\]
\begin{itemize}[leftmargin=1.2em]
  \item[\color{cTeal}\faUser] \textbf{edad}: $e^{0,0199}=1,020$. Cada año adicional multiplica los odds de
  enfermedad por 1,02 (+2\,\%); en 10 años, $e^{10(0,0199)}=1,22$ (+22\,\%), con las demás variables fijas.
  No es significativo ($p=0,31$).
  \item[\color{cRed}\faHeartbeat] \textbf{angina\_ejercicio}: $e^{1,325}=3,76$. Un paciente con angina
  inducida por ejercicio tiene odds de enfermedad 3,8 veces mayores que uno sin angina, a igualdad del resto.
  \item[\color{cRed}\faChartLine] \textbf{depresion\_st}: $e^{0,583}=1,79$. Cada milímetro adicional de
  depresión del segmento ST multiplica los odds por 1,79 (+79\,\%).
\end{itemize}

\pregunta{d)}{Odds ratios e IC 95\,\%}{Calcule e interprete los odds ratios y sus intervalos de confianza al
95\,\%.}
\codigo{log-or}
\salida{log-or}
Los intervalos son por verosimilitud perfil (\var{confint} de un \var{glm}), mejores que $e^{\hat\beta\pm1,96\,EE}$
cuando la log-verosimilitud no es simétrica.
\begin{center}\small
\begin{tabular}{lrrl}
\toprule
Variable & OR & IC 95\,\% & Lectura \\ \midrule
sexo (hombre) & 5,22 & [2,62 ; 10,93] & odds 5 veces mayores en hombres \\
angina\_ejercicio (sí) & 3,76 & [1,96 ; 7,38] & odds casi 4 veces mayores \\
depresion\_st & 1,79 & [1,35 ; 2,42] & +79\,\% por mm \\
colesterol & 1,006 & [1,0003 ; 1,012] & +6,4\,\% cada 10 mg/dl \\
frecuencia\_cardiaca\_max & 0,971 & [0,955 ; 0,986] & $-26$\,\% cada 10 lpm ($0,971^{10}=0,74$) \\
presion\_reposo & 1,014 & [0,997 ; 1,033] & incluye 1: no concluyente \\
edad & 1,020 & [0,982 ; 1,060] & incluye 1: no concluyente \\ \bottomrule
\end{tabular}
\end{center}

\res{Los IC que no contienen al 1 coinciden con las variables significativas. Los factores de riesgo más
fuertes son ser hombre y tener angina de esfuerzo; una frecuencia cardíaca máxima alta es protectora.}

\pregunta{e)}{Calidad del ajuste: razón de verosimilitudes y matriz de confusión}{Evalúe la calidad del ajuste
utilizando la razón de verosimilitudes, la matriz de confusión y la tasa de clasificación correcta (umbral
0,5); interprete exactitud, sensibilidad y especificidad en el contexto del diagnóstico.}
\codigo{log-lr}
\salida{log-lr}
\[
G^2=-2\big[\ln L_0-\ln L_1\big]=409,95-274,16=135,8\;\sim\;\chi^2_7,\qquad p<2\times10^{-16}
\]
Rechazo $H_0:\beta_1=\dots=\beta_7=0$: el modelo explica significativamente más que solo la prevalencia.
El $R^2$ de McFadden (0,33) está en el rango que se considera un buen ajuste (0,2--0,4).
\codigo{log-confusion}
\salida{log-confusion}
\begin{center}\small
\begin{tabular}{lcc}
\toprule
& \multicolumn{2}{c}{Clase predicha} \\ \cmidrule(lr){2-3}
Clase real & No & Sí \\ \midrule
No & VN $=133$ & FP $=27$ \\
Sí & FN $=37$ & VP $=100$ \\ \bottomrule
\end{tabular}
\end{center}
\[
\text{Exactitud}=\frac{100+133}{297}=0,785
\]
\[
\text{Sensibilidad}=\frac{100}{100+37}=0,730,\qquad
\text{Especificidad}=\frac{133}{133+27}=0,831
\]
\begin{itemize}[leftmargin=1.2em]
  \item[\color{cTeal}\faCheck] \textbf{Exactitud 78,5\,\%}: clasifica bien a casi 4 de cada 5 pacientes
  (la regla trivial ``todos sanos'' acertaría 53,9\,\%).
  \item[\color{cRed}\faExclamationTriangle] \textbf{Sensibilidad 73,0\,\%}: detecta 100 de 137 enfermos;
  \textbf{37 enfermos quedan como sanos}. En diagnóstico cardíaco este es el error caro, porque un enfermo
  no detectado no recibe tratamiento.
  \item[\color{cTeal}\faIcon{shield-alt}] \textbf{Especificidad 83,1\,\%}: 27 sanos serían derivados a pruebas
  adicionales sin necesidad; es un costo, pero menor.
\end{itemize}

\res{El modelo ajusta significativamente ($G^2=135,8$) y clasifica bien el 78,5\,\%. Con umbral 0,5 es más
específico que sensible; para tamizaje convendría bajar el umbral (ver pregunta f).}

\pregunta{f)}{Curva ROC y AUC}{Analice la capacidad predictiva del modelo mediante la curva ROC e interprete el
área bajo la curva (AUC).}
\figcod{log-roc}{Curva ROC con el punto de Youden (arriba) y probabilidades estimadas por clase (abajo).}
\codigo{log-roc}
\salida{log-roc}
$\text{AUC}=0,858$ (IC 95\,\% DeLong: 0,816--0,900): si tomo al azar un enfermo y un sano, en el 85,8\,\% de
los casos el modelo le asigna mayor probabilidad al enfermo. Es una discriminación \emph{buena} (0,8--0,9) y el
IC está lejos de 0,5. El umbral de Youden (0,443) sube la sensibilidad a 78,1\,\% con especificidad de 80\,\%,
un mejor equilibrio para tamizaje que 0,5.

\res{AUC $=0,858$: el modelo discrimina bien entre enfermos y sanos.}

\pregunta{g)}{Limitaciones y extensiones}{Discuta las limitaciones del modelo y proponga posibles extensiones.}
Todas las métricas anteriores se calcularon sobre los mismos 297 pacientes del ajuste, lo que las hace
optimistas. Lo verifiqué con validación cruzada de 10 pliegues:
\codigo{log-cv}
\salida{log-cv}
Fuera de muestra el AUC baja de 0,858 a 0,836 y la exactitud de 78,5\,\% a 76,1\,\%: hay algo de
sobreajuste, pero el modelo generaliza razonablemente.

\noindent\textbf{Limitaciones}
\begin{itemize}[leftmargin=1.2em]
  \item[\color{cRed}\faMinusCircle] Solo 297 pacientes de un único hospital (Cleveland): poca validez externa.
  \item[\color{cRed}\faMinusCircle] Supone que el logit es lineal en cada variable continua (edad, colesterol,
  frecuencia cardíaca, depresión ST) y no incluye interacciones (por ejemplo, sexo $\times$ edad).
  \item[\color{cRed}\faMinusCircle] Deja fuera predictores muy informativos de la base: tipo de dolor torácico,
  vasos coloreados por fluoroscopía, talasemia y pendiente ST.
  \item[\color{cRed}\faMinusCircle] Es un estudio transversal: los OR son asociaciones, no efectos causales.
\end{itemize}
\noindent\textbf{Extensiones}
\begin{itemize}[leftmargin=1.2em]
  \item[\color{cTeal}\faPlusCircle] Agregar las variables omitidas y seleccionar con AIC/BIC o con
  regularización (lasso).
  \item[\color{cTeal}\faPlusCircle] Términos no lineales (splines en edad y frecuencia cardíaca) o un GAM.
  \item[\color{cTeal}\faPlusCircle] Usar \var{diagnostico_original} (0--4) en un modelo logístico ordinal de
  severidad.
  \item[\color{cTeal}\faPlusCircle] Elegir el umbral según el costo relativo FN/FP y validar en una cohorte externa.
\end{itemize}

\end{document}
